\chapter{C 题决策树与大模拟方法}\label{ch:07}

\section{C 题定位}\label{sec:07-1}

C 题的长题面常见多个对象、几种命令和一些特殊规则。读题和组织代码会占不少时间，先把这些内容分清楚，后面才容易查错。

\begin{itemize}
  \item 建议目标时间：60--100 分钟。
  \item 目标：能满分则满分；若细节太多，先完成题目明确给出的子任务。
  \item 高频主题：大模拟、字符串解析、状态维护、网格搜索、图论基础、数据结构基础、简单 DP。
  \item 容易出错的地方：漏规则、状态没记全、不同命令混在同一段代码里。
\end{itemize}

可以先列一张对象和命令表，确认每种操作会改哪些字段。

\section{C 题开题流程}\label{sec:07-2}

\begin{enumerate}
  \item 快速浏览题面，判断是大模拟还是算法题。
  \item 圈出所有对象：人、文件、任务、节点、格子、命令、事件。
  \item 圈出所有状态：位置、时间、分数、开关、权限、颜色、方向。
  \item 圈出所有操作：新增、删除、查询、移动、更新、比较。
  \item 圈出优先级和并列规则。
  \item 设计结构体和函数。
  \item 先实现最基础规则，通过样例的一部分。
  \item 再逐条补特殊规则。
\end{enumerate}

\section{C 题读题三遍法}\label{sec:07-3}

题面较长时可以分几遍读：先看要模拟什么，再记状态和命令，最后核对特殊规则。

\subsection{第一遍：只判断题型}

第一遍快速读，不抠细节，只回答：

\begin{itemize}
  \item 是大模拟，还是算法题？
  \item 输入是一次性数据，还是很多命令？
  \item 输出是每次查询输出，还是最后统一输出？
  \item 数据范围是否明显要求优化？
\end{itemize}

第一遍还没理清时，可以结合样例看看一次操作前后发生了什么。

\subsection{第二遍：圈对象、状态、操作}

第二遍要动笔。把题目中的名词和动词分开。

\begin{longtable}{p{3cm}p{5cm}p{6cm}}
\toprule
类别 & 怎么找 & 例子 \\
\midrule
\endfirsthead
\toprule
类别 & 怎么找 & 例子 \\
\midrule
\endhead
对象 & 被操作的名词 & 用户、文件、进程、任务、车辆、节点、格子。 \\
属性 & 对象身上的数据 & 编号、名字、位置、时间、状态、得分。 \\
操作 & 题目中的动词或命令 & 新增、删除、移动、查询、合并、修改。 \\
规则 & 操作发生时的限制 & 如果不存在则忽略；如果冲突按优先级。 \\
输出 & 什么时候打印 & 每次 query、最后统计、发生错误时。 \\
\bottomrule
\end{longtable}

\subsection{第三遍：找坑点}

第三遍只找容易错的细节：

\begin{itemize}
  \item 相同时间、相同分数、相同优先级怎么处理。
  \item 编号从 0 还是 1 开始。
  \item 找不到对象时怎么处理。
  \item 重复添加时怎么处理。
  \item 删除后是否还能查询。
  \item 最后一条命令或最后一个事件是否要特殊处理。
  \item 输出顺序按输入顺序、编号顺序、字典序还是时间顺序。
\end{itemize}

\section{从题面到代码的转换}\label{sec:07-4}

\subsection{对象变成 struct}

题目中的每类对象，优先考虑一个结构体。

\begin{lstlisting}
struct Task {
    int id;
    int startTime;
    int endTime;
    int priority;
    int status;
};
\end{lstlisting}

\subsection{编号查找用 map 或 vector}

如果编号是 $1$ 到 $n$：

\begin{lstlisting}
vector<Task> tasks(n + 1);
\end{lstlisting}

如果编号很大或是字符串：

\begin{lstlisting}
map<string, Task> tasks;
unordered_map<int, Task> tasksById;
\end{lstlisting}

\subsection{命令变成函数}

每个命令单独写函数：

\begin{lstlisting}
void addTask(int id, int priority) {
    // add or update
}

void deleteTask(int id) {
    // delete if exists
}

void queryTask(int id) {
    // print answer
}
\end{lstlisting}

\subsection{规则变成 if}

题目中的条件要落实到代码里。例如，查询前先判断对象是否存在，再判断它的状态：

\begin{lstlisting}
if (!tasks.count(id)) {
    cout << "Not Found\n";
    return;
}

if (tasks[id].status == FINISHED) {
    cout << "Finished\n";
    return;
}
\end{lstlisting}

把几条不同规则分开写，样例不对时比较容易找到是哪一步。

\section{状态维护和容器选择}\label{sec:07-5}

\subsection{容器选择表}

\begin{longtable}{p{4cm}p{5cm}p{5cm}}
\toprule
需求 & 推荐容器 & 原因 \\
\midrule
\endfirsthead
\toprule
需求 & 推荐容器 & 原因 \\
\midrule
\endhead
编号是 $1$ 到 $n$ & \Code{vector<Obj>} & 直接下标访问，最快。 \\
编号很大但唯一 & \Code{unordered\_map<int,Obj>} & 平均快速查找。 \\
编号是字符串 & \Code{map<string,Obj>} 或 \Code{unordered\_map} & 字符串到对象映射。 \\
需要按 key 有序输出 & \Code{map} & 自动按 key 排序。 \\
需要判断是否存在 & \Code{set} / \Code{unordered\_set} & 查重。 \\
需要动态最小/最大 & \Code{set} / \Code{priority\_queue} & 维护极值。 \\
需要前驱后继 & \Code{set} & \Code{lower\_bound}。 \\
需要按进入顺序处理 & \Code{queue} & FIFO。 \\
需要最近使用顺序 & \Code{list + map} & LRU。 \\
需要父子层级 & \Code{vector<Node>} + child & 树结构。 \\
\bottomrule
\end{longtable}

\subsection{操作复杂度估算}

\begin{longtable}{p{4cm}p{4cm}p{6cm}}
\toprule
容器 & 常见复杂度 & 注意 \\
\midrule
\endfirsthead
\toprule
容器 & 常见复杂度 & 注意 \\
\midrule
\endhead
\Code{vector} 下标访问 & $O(1)$ & 中间插入删除慢。 \\
\Code{map} 查找插入 & $O(\log n)$ & 有序但常数较大。 \\
\Code{unordered\_map} 查找插入 & 平均 $O(1)$ & 无序，最坏情况不稳定。 \\
\Code{set} 查找插入删除 & $O(\log n)$ & 可找前驱后继。 \\
\Code{priority\_queue} 取最优 & $O(\log n)$ 插入删除堆顶 & 不能直接删除中间元素。 \\
\Code{queue} 进出队 & $O(1)$ & 只能访问队首。 \\
\Code{list} 插入删除 & 已有迭代器时 $O(1)$ & 不能随机访问。 \\
\bottomrule
\end{longtable}

\subsection{C 题容器选择错误信号}

\begin{itemize}
  \item 每次命令都全表扫描，且命令数很大：应该用 map/set/hash。
  \item 每次取最大值都排序：应该用 set 或堆。
  \item 每次删除 vector 中间元素：可能需要标记删除、list 或 set。
  \item 需要有序输出却用了 unordered\_map：输出顺序会错。
  \item priority\_queue 里元素更新了但旧值还在：需要懒删除。
\end{itemize}

\section{C 题总决策树}\label{sec:07-6}

\begin{longtable}{p{4cm}p{5cm}p{5cm}}
\toprule
题面信号 & 优先怀疑 & 首选方法 \\
\midrule
\endfirsthead
\toprule
题面信号 & 优先怀疑 & 首选方法 \\
\midrule
\endhead
题面很长，规则很多 & 大模拟 & 结构体 + 函数拆解 \\
输入是一串命令 & 命令解析模拟 & stringstream / split / 状态表 \\
按时间发生事件 & 事件模拟 & 事件排序 / 优先队列 \\
对象有多种状态 & 状态机模拟 & enum/int 表示状态 \\
文件路径、配置、表达式 & 字符串解析 & 逐字符扫描 / 栈 \\
网格、地图、移动 & 网格模拟或 BFS/DFS & 方向数组、vis、dist \\
图、连通、最短 & 图论 & BFS / Dijkstra / 并查集 \\
区间修改查询 & 数据结构 & 前缀和 / 差分 / 树状数组 / 线段树 \\
求最优值且有阶段 & DP & 先定义状态 \\
\bottomrule
\end{longtable}

\section{C 题常见题型}\label{sec:07-7}

这张表用于快速判断 C 题该从哪个方向切入。

\begin{longtable}{p{4cm}p{5cm}p{5cm}}
\toprule
题目长相 & 本质 & 先翻哪里 \\
\midrule
\endfirsthead
\toprule
题目长相 & 本质 & 先翻哪里 \\
\midrule
\endhead
很多命令，每条命令改变系统状态 & 命令模拟 & 命令解析、状态设计、函数化 \\
按时间发生、开始结束、排队 & 事件模拟 & 事件排序、优先级、队列模拟 \\
字符串里有路径、参数、配置 & 字符串解析 & split、key=value、路径规范化 \\
输入像 URL、路由、参数 & 路由匹配 & URL 路由和参数解析 \\
输入像 JSON、配置、嵌套括号 & 结构解析 & 递归解析、栈、逐字符扫描 \\
输入像 Markdown、HTML、轻量标记 & 标记语言模拟 & 行扫描、状态切换 \\
有用户、角色、权限、规则 & 规则匹配 & 权限和规则匹配 \\
有表格、矩阵、坐标变换 & 表格模拟 & 矩阵操作、行列映射 \\
有缓存、队列、窗口、淘汰 & 数据结构模拟 & 队列、set、map、懒删除 \\
有目录、树、层级关系 & 树结构模拟 & 目录树、父子关系、DFS \\
有地图并带钥匙/状态 & 状态搜索 & 带状态 BFS \\
\bottomrule
\end{longtable}

\section{C 题细分决策树}\label{sec:07-8}

\subsection{第一问：这是大模拟还是算法题}

\begin{longtable}{p{5cm}p{5cm}p{4cm}}
\toprule
题面特征 & 判断 & 下一步 \\
\midrule
\endfirsthead
\toprule
题面特征 & 判断 & 下一步 \\
\midrule
\endhead
规则很多、命令很多、对象很多 & 大模拟 & 拆对象、状态、命令 \\
题面不长但数据范围大 & 算法题 & 回到 B/D 算法决策树 \\
有地图和最短步数 & 搜索题 & BFS / 状态 BFS \\
有依赖关系 & 图论题 & 拓扑 / 最短路 / 并查集 \\
有复杂字符串格式 & 解析题 & 字符串解析决策树 \\
有大量修改查询 & 数据结构题 & 前缀/差分/树状数组/线段树 \\
\bottomrule
\end{longtable}

\subsection{大模拟细分树}

\begin{longtable}{p{5cm}p{5cm}p{4cm}}
\toprule
题面信号 & 类型 & 首选结构 \\
\midrule
\endfirsthead
\toprule
题面信号 & 类型 & 首选结构 \\
\midrule
\endhead
每行一个操作名 & 命令系统 & \Code{handleCommand} 分发 \\
对象有生命周期 & 增删改查模拟 & \Code{map<int,Obj>} 或 \Code{vector<Obj>} \\
对象有多个状态 & 状态机 & 状态常量 + 转移函数 \\
按时间发生 & 事件模拟 & \Code{Event} + 排序 \\
每次取最优对象 & 优先级模拟 & 堆 / set / 懒删除堆 \\
行列、表格、棋盘 & 表格模拟 & 二维数组 / 行列映射 \\
路径、目录、父子层级 & 树结构模拟 & 节点数组 + child map \\
规则匹配、权限判断 & 规则系统 & \Code{Rule} + \Code{match} \\
\bottomrule
\end{longtable}

\subsection{字符串解析细分树}

\begin{longtable}{p{5cm}p{5cm}p{4cm}}
\toprule
字符串形态 & 优先方法 & 小心 \\
\midrule
\endfirsthead
\toprule
字符串形态 & 优先方法 & 小心 \\
\midrule
\endhead
按空格分隔 & \Code{stringstream} & 前面是否要 \Code{ignore} \\
按固定字符分隔 & \Code{split} & 连续分隔符 \\
\Code{key=value} & \Code{find + substr} & 值是否含空格 \\
路径 \Code{/a/b/c} & 路径 split / normalize & 根目录、\Code{..} \\
括号匹配 & 栈 & 空栈和最后剩余 \\
表达式求值 & 扫描 / 递归解析 & 优先级、括号 \\
URL 参数 & 先按 \Code{?}，再按 \Code{\&} & 参数为空 \\
嵌套 JSON 类 & 逐字符扫描 / 栈 & 先利用题目限制 \\
\bottomrule
\end{longtable}

\subsection{图网格细分树}

\begin{longtable}{p{5cm}p{5cm}p{4cm}}
\toprule
题面信号 & 类型 & 首选模板 \\
\midrule
\endfirsthead
\toprule
题面信号 & 类型 & 首选模板 \\
\midrule
\endhead
最少步数、每步代价相同 & 无权最短路 & BFS \\
多个起点最近距离 & 多源 BFS & 所有起点入队 \\
有钥匙、方向、模式等额外状态 & 状态 BFS & \Code{dist[x][y][state]} \\
统计区域大小 & 连通块 & DFS / BFS \\
正权边最短路 & 正权图 & Dijkstra \\
只问是否连通 & 连通性 & 并查集 / DFS \\
任务依赖顺序 & DAG & 拓扑排序 \\
\bottomrule
\end{longtable}

\section{大模拟总方法}\label{sec:07-9}

\subsection{大模拟什么题会用到}

\begin{itemize}
  \item 题面长度明显长。
  \item 有很多“如果……否则……”规则。
  \item 输入包含多条命令。
  \item 有多个对象，每个对象有状态。
  \item 输出要求根据最终状态或每次查询得到。
\end{itemize}

\subsection{大模拟代码骨架}

\begin{lstlisting}
#include <bits/stdc++.h>
using namespace std;
using ll = long long;

struct Item {
    int id;
    string name;
    int state;
};

int n, m;
vector<Item> items;

void readInput() {
    cin >> n >> m;
    items.resize(n);
    for (int i = 0; i < n; i++) {
        cin >> items[i].id >> items[i].name;
        items[i].state = 0;
    }
}

void handleCommand(const string& cmd) {
    if (cmd == "ADD") {
        // handle add
    } else if (cmd == "DEL") {
        // handle delete
    } else if (cmd == "QUERY") {
        // handle query
    }
}

void solve() {
    for (int i = 0; i < m; i++) {
        string cmd;
        cin >> cmd;
        handleCommand(cmd);
    }
}

int main() {
    ios::sync_with_stdio(false);
    cin.tie(nullptr);

    readInput();
    solve();
    return 0;
}
\end{lstlisting}

\subsection{为什么要拆函数}

\begin{itemize}
  \item 读入、处理、输出分开，查错更容易。
  \item 每个命令一个函数，避免主函数过长。
  \item 特殊规则可以单独补，不容易破坏已有逻辑。
  \item 样例错时能定位是哪条规则错。
\end{itemize}

\section{长题面拆解表}\label{sec:07-10}

读 C 题时，在草稿纸上画这个表：

\begin{longtable}{p{3cm}p{5cm}p{6cm}}
\toprule
项目 & 要记录什么 & 示例 \\
\midrule
\endfirsthead
\toprule
项目 & 要记录什么 & 示例 \\
\midrule
\endhead
对象 & 题目中被操作的东西 & 用户、文件、任务、格子、节点。 \\
状态 & 对象身上的变量 & 位置、分数、权限、颜色、是否激活。 \\
命令 & 输入中每种操作 & add、delete、query、move。 \\
规则 & 每条命令如何改变状态 & 如果存在则更新，否则新增。 \\
优先级 & 相同条件时谁先 & 时间小优先、编号小优先。 \\
边界 & 特殊或非法情况 & 找不到、重复、越界、空集合。 \\
输出 & 什么时候输出什么 & 每次 query 输出，最后统一输出。 \\
\bottomrule
\end{longtable}

\section{命令解析}\label{sec:07-11}

\subsection{固定格式命令}

如果每条命令格式固定：

\begin{lstlisting}
string op;
cin >> op;
if (op == "add") {
    int id, x;
    cin >> id >> x;
} else if (op == "query") {
    int id;
    cin >> id;
}
\end{lstlisting}

\subsection{整行命令}

如果命令中可能有空格，先读整行。

\begin{lstlisting}
int q;
cin >> q;
cin.ignore(numeric_limits<streamsize>::max(), '\n');

while (q--) {
    string line;
    getline(cin, line);
    stringstream ss(line);
    string op;
    ss >> op;

    if (op == "add") {
        string name;
        int x;
        ss >> name >> x;
    }
}
\end{lstlisting}

\subsection{手写 split}

按指定字符分割路径、日期、配置项：

\begin{lstlisting}
vector<string> split(const string& s, char sep) {
    vector<string> res;
    string cur;
    for (char c : s) {
        if (c == sep) {
            res.push_back(cur);
            cur.clear();
        } else {
            cur.push_back(c);
        }
    }
    res.push_back(cur);
    return res;
}
\end{lstlisting}

\subsection{命令分发模板}

命令较多时，可以让每种命令对应一个函数，主循环只负责读入和调用。

\begin{lstlisting}
void handleAdd() {
    int id, x;
    cin >> id >> x;
    // add logic
}

void handleDelete() {
    int id;
    cin >> id;
    // delete logic
}

void handleQuery() {
    int id;
    cin >> id;
    // query logic
}

void handleCommand() {
    string op;
    cin >> op;
    if (op == "add") handleAdd();
    else if (op == "delete") handleDelete();
    else if (op == "query") handleQuery();
}
\end{lstlisting}

\subsection{命令解析易错点}

\begin{itemize}
  \item 命令参数个数是否固定。
  \item 参数中是否可能包含空格。
  \item 命令大小写是否敏感。
  \item 非法命令是否会出现。
  \item 每条命令是否都可能输出。
  \item 查询不存在对象时输出什么。
\end{itemize}

\section{状态设计}\label{sec:07-12}

\subsection{用结构体保存对象}

\begin{lstlisting}
struct User {
    int id;
    string name;
    int score;
    bool active;
};

map<int, User> users;
\end{lstlisting}

\subsection{状态编号}

状态不多时可以用整数常量：

\begin{lstlisting}
const int WAITING = 0;
const int RUNNING = 1;
const int FINISHED = 2;
const int FAILED = 3;
\end{lstlisting}

也可以用 \Code{enum}：

\begin{lstlisting}
enum Status {
    WAITING,
    RUNNING,
    FINISHED,
    FAILED
};
\end{lstlisting}

\subsection{状态转移表}

复杂状态题建议先画表：

\begin{longtable}{p{3cm}p{3cm}p{4cm}p{4cm}}
\toprule
当前状态 & 操作 & 新状态 & 输出/副作用 \\
\midrule
\endfirsthead
\toprule
当前状态 & 操作 & 新状态 & 输出/副作用 \\
\midrule
\endhead
WAITING & start & RUNNING & 记录开始时间 \\
RUNNING & finish & FINISHED & 计算得分 \\
RUNNING & fail & FAILED & 输出错误信息 \\
\bottomrule
\end{longtable}

\section{状态机模拟}\label{sec:07-13}

\subsection{什么题会用到}

题面出现：

\begin{itemize}
  \item 一个对象有多个状态。
  \item 不同状态下，同一个命令效果不同。
  \item 状态按规则切换。
  \item 非法状态转换需要忽略或报错。
\end{itemize}

\subsection{状态机模板}

\begin{lstlisting}
const int WAITING = 0;
const int RUNNING = 1;
const int FINISHED = 2;

struct Job {
    int id;
    int status;
    int score;
};

void start(Job& job) {
    if (job.status != WAITING) {
        return;
    }
    job.status = RUNNING;
}

void finish(Job& job) {
    if (job.status != RUNNING) {
        return;
    }
    job.status = FINISHED;
    job.score += 10;
}
\end{lstlisting}

\subsection{状态机检查}

\begin{itemize}
  \item 初始状态是什么。
  \item 每个操作允许在哪些状态下发生。
  \item 不允许的操作是忽略、报错，还是改变成特殊状态。
  \item 终止状态是否还能被修改。
  \item 状态改变时是否还要更新其他字段。
\end{itemize}

\section{事件模拟}\label{sec:07-14}

\subsection{什么题会用到}

题面出现：

\begin{itemize}
  \item 按时间顺序发生事件。
  \item 事件有开始、结束、到达、离开。
  \item 同一时间有多个事件。
  \item 需要维护当前正在发生的对象集合。
\end{itemize}

\subsection{事件结构体}

\begin{lstlisting}
struct Event {
    int time;
    int type;
    int id;
};

vector<Event> events;
\end{lstlisting}

\subsection{事件排序}

\begin{lstlisting}
sort(events.begin(), events.end(), [](const Event& a, const Event& b) {
    if (a.time != b.time) return a.time < b.time;
    return a.type < b.type; // depends on statement
});
\end{lstlisting}

\subsection{事件处理模板}

\begin{lstlisting}
for (const Event& e : events) {
    if (e.type == 0) {
        // start event
    } else if (e.type == 1) {
        // end event
    } else {
        // query event
    }
}
\end{lstlisting}

\subsection{同一时间事件顺序}

这是事件模拟最容易错的地方。一定要从题面确定：

\begin{itemize}
  \item 同一时间，结束事件先处理还是开始事件先处理。
  \item 查询是在事件前还是事件后。
  \item 如果没有说明，要从样例推断，但不要过度猜。
  \item 排序比较器中必须体现这个顺序。
\end{itemize}

\section{优先级模拟}\label{sec:07-15}

\subsection{什么题会用到}

题面出现：

\begin{itemize}
  \item 每次选择优先级最高的任务。
  \item 优先级相同按时间、编号、字典序。
  \item 动态加入任务。
  \item 当前最大、当前最小。
\end{itemize}

\subsection{priority\_queue 模板}

\begin{lstlisting}
struct Task {
    int priority;
    int time;
    int id;
};

struct Cmp {
    bool operator()(const Task& a, const Task& b) const {
        if (a.priority != b.priority) return a.priority < b.priority;
        if (a.time != b.time) return a.time > b.time;
        return a.id > b.id;
    }
};

priority_queue<Task, vector<Task>, Cmp> pq;
\end{lstlisting}

这个比较器表示：

\begin{itemize}
  \item 优先级大的先出。
  \item 优先级相同，时间小的先出。
  \item 时间相同，编号小的先出。
\end{itemize}

\Warn{priority\_queue 的比较器和 sort 相反，容易写错。}不确定时，用几个任务手动推一下 \Code{top()} 是谁。

\section{复杂排序和排名}\label{sec:07-16}

\subsection{什么题会用到}

题面出现：

\begin{itemize}
  \item 多个排序关键字。
  \item 排名、积分榜、成绩表。
  \item 相同分数按编号、时间、字典序排序。
  \item 并列排名或去重排名。
\end{itemize}

\subsection{复杂排序模板}

\begin{lstlisting}
struct Player {
    int id;
    int score;
    int penalty;
    string name;
};

sort(a.begin(), a.end(), [](const Player& x, const Player& y) {
    if (x.score != y.score) return x.score > y.score;
    if (x.penalty != y.penalty) return x.penalty < y.penalty;
    if (x.name != y.name) return x.name < y.name;
    return x.id < y.id;
});
\end{lstlisting}

\subsection{并列排名}

\begin{lstlisting}
int rank = 1;
for (int i = 0; i < (int)a.size(); i++) {
    if (i > 0 && a[i].score != a[i - 1].score) {
        rank = i + 1;
    }
    cout << a[i].id << ' ' << rank << '\n';
}
\end{lstlisting}

\subsection{复杂排序检查}

\begin{itemize}
  \item 每个关键字是升序还是降序。
  \item 相同关键字时是否继续比较。
  \item 是否要求稳定排序。
  \item 字符串按字典序还是输入顺序。
  \item 比较器相等时不能返回 \Code{true}。
\end{itemize}

\section{表格和矩阵模拟}\label{sec:07-17}

\subsection{什么题会用到}

题面出现：

\begin{itemize}
  \item 表格、座位、棋盘、网格。
  \item 行列操作。
  \item 旋转、翻转、交换行列。
  \item 查询某个单元格。
\end{itemize}

\subsection{矩阵读入和输出}

\begin{lstlisting}
int n, m;
cin >> n >> m;
vector<vector<int>> a(n, vector<int>(m));
for (int i = 0; i < n; i++) {
    for (int j = 0; j < m; j++) {
        cin >> a[i][j];
    }
}
\end{lstlisting}

\subsection{交换两行}

\begin{lstlisting}
int r1, r2;
cin >> r1 >> r2;
--r1; --r2;
swap(a[r1], a[r2]);
\end{lstlisting}

\subsection{交换两列}

\begin{lstlisting}
int c1, c2;
cin >> c1 >> c2;
--c1; --c2;
for (int i = 0; i < n; i++) {
    swap(a[i][c1], a[i][c2]);
}
\end{lstlisting}

\subsection{顺时针旋转矩阵}

\begin{lstlisting}
vector<vector<int>> rotateClockwise(const vector<vector<int>>& a) {
    int n = (int)a.size();
    if (n == 0) return {};
    int m = (int)a[0].size();
    vector<vector<int>> b(m, vector<int>(n));
    for (int i = 0; i < n; i++) {
        for (int j = 0; j < m; j++) {
            b[j][n - 1 - i] = a[i][j];
        }
    }
    return b;
}
\end{lstlisting}

\subsection{表格模拟易错点}

\begin{itemize}
  \item 题目行列通常从 1 开始，代码常从 0 开始。
  \item 旋转后行数和列数会交换。
  \item 如果操作很多，真实移动整个矩阵可能超时，要考虑只维护行列映射。
  \item 输出每行空格格式。
\end{itemize}

\section{目录树和路径模拟}\label{sec:07-18}

\subsection{什么题会用到}

题面出现：

\begin{itemize}
  \item 文件、目录、路径。
  \item \Code{/a/b/c}、相对路径、绝对路径。
  \item 创建、删除、查询文件。
  \item 父目录、子目录。
\end{itemize}

\subsection{路径规范化}

下面处理绝对路径的 \Code{.} 和 \Code{..}；相对路径需先拼接当前目录。依赖前文 split：

\begin{lstlisting}
vector<string> normalizePath(const string& path) {
    vector<string> st;
    vector<string> parts = split(path, '/');
    for (string part : parts) {
        if (part.empty() || part == ".") {
            continue;
        } else if (part == "..") {
            if (!st.empty()) st.pop_back();
        } else {
            st.push_back(part);
        }
    }
    return st;
}
\end{lstlisting}

\subsection{路径转字符串}

\begin{lstlisting}
string joinPath(const vector<string>& parts) {
    if (parts.empty()) return "/";
    string res;
    for (const string& p : parts) {
        res += "/";
        res += p;
    }
    return res;
}
\end{lstlisting}

\subsection{目录树节点}

\begin{lstlisting}
struct Node {
    string name;
    map<string, int> child;
    bool isFile = false;
};

vector<Node> tree;
\end{lstlisting}

新建节点：

\begin{lstlisting}
int newNode(const string& name, bool isFile) {
    Node node;
    node.name = name;
    node.isFile = isFile;
    tree.push_back(node);
    return (int)tree.size() - 1;
}
\end{lstlisting}

\subsection{路径题易错点}

\begin{itemize}
  \item 根目录怎么表示。
  \item 连续斜杠是否可能出现。
  \item 路径末尾斜杠是否影响结果。
  \item 查询不存在路径时输出什么。
  \item 文件和目录是否允许同名。
\end{itemize}

\section{URL 路由和参数解析}\label{sec:07-19}

\subsection{什么题会用到}

题面出现：

\begin{itemize}
  \item URL、路径、路由。
  \item \Code{/user/123/profile}。
  \item 查询参数，例如 \Code{?a=1\&b=2}。
  \item 通配符、变量匹配。
\end{itemize}

\subsection{拆分 URL}

\begin{lstlisting}
struct Url {
    string path;
    map<string, string> query;
};

Url parseUrl(const string& s) {
    Url res;
    size_t qpos = s.find('?');
    if (qpos == string::npos) {
        res.path = s;
        return res;
    }

    res.path = s.substr(0, qpos);
    string qs = s.substr(qpos + 1);
    vector<string> pairs = split(qs, '&');
    for (string p : pairs) {
        if (p.empty()) continue;
        size_t eq = p.find('=');
        if (eq == string::npos) {
            res.query[p] = "";
        } else {
            string key = p.substr(0, eq);
            string val = p.substr(eq + 1);
            res.query[key] = val;
        }
    }
    return res;
}
\end{lstlisting}

\subsection{简单路由匹配}

模式 \Code{/user/<id>} 匹配 \Code{/user/123}。先定义后文的 \Code{parsePath}，或在函数前声明它。此简化版忽略连续/末尾斜杠，仅支持整段参数，不支持 URL 映射真题的类型参数和剩余路径规则：

\begin{lstlisting}
bool matchRoute(const string& pattern, const string& path,
                map<string, string>& params) {
    params.clear();
    map<string, string> candidate;
    vector<string> a = parsePath(pattern);
    vector<string> b = parsePath(path);
    if (a.size() != b.size()) return false;

    for (int i = 0; i < (int)a.size(); i++) {
        if (!a[i].empty() && a[i][0] == '<' && a[i].back() == '>') {
            string name = a[i].substr(1, (int)a[i].size() - 2);
            candidate[name] = b[i];
        } else if (a[i] != b[i]) {
            return false;
        }
    }
    params = move(candidate);
    return true;
}
\end{lstlisting}

\subsection{URL 题易错点}

\begin{itemize}
  \item 有没有查询参数。
  \item 参数值是否可能为空。
  \item 路径末尾斜杠是否等价。
  \item 多个路由都匹配时，第一条生效还是更具体的生效。
  \item URL 是否需要解码，题目不要求就不要自己加。
\end{itemize}

\section{配置和 JSON 类解析入口}\label{sec:07-20}

\subsection{什么题会用到}

题面出现：

\begin{itemize}
  \item 键值对。
  \item 花括号、方括号、冒号、逗号。
  \item 嵌套对象。
  \item 查询某个字段。
\end{itemize}

\subsection{简单键值配置}

每行形如 \Code{key=value}：

\begin{lstlisting}
int n;
cin >> n;
map<string, string> conf;
for (int i = 0; i < n; i++) {
    string line;
    cin >> line;
    int pos = (int)line.find('=');
    string key = line.substr(0, pos);
    string value = line.substr(pos + 1);
    conf[key] = value;
}
\end{lstlisting}

\subsection{读取带空格的配置值}

\begin{lstlisting}
string line;
getline(cin, line);
int pos = (int)line.find('=');
string key = line.substr(0, pos);
string value = line.substr(pos + 1);
\end{lstlisting}

\subsection{嵌套结构处理思路}

JSON 类题先看题目允许哪些数据类型和格式，只实现题目实际需要的部分：

\begin{itemize}
  \item 是否只有字符串和对象，没有数组。
  \item 是否保证格式合法。
  \item 查询是否只按点号路径。
  \item 是否只需要保存完整路径到值的映射。
\end{itemize}

如果只需要查询路径，可以把嵌套字段拍平成 key：

\begin{lstlisting}
// Example flattened keys:
// user.name -> Alice
// user.age -> 18
map<string, string> value;
\end{lstlisting}

\Warn{配置/JSON 题先利用限制。}CSP 题通常不会要求你写工业级 JSON 解析器，重点是按题目给定范围实现。

\section{字符串解析}\label{sec:07-21}

\subsection{什么题会用到}

遇到下面这些格式时，需要先想清楚字符串怎样拆开：

\begin{itemize}
  \item 路径，例如 \Code{/a/b/c}。
  \item 配置项，例如 \Code{key=value}。
  \item 表达式，例如括号、加减乘除。
  \item 带引号的字符串。
  \item 嵌套结构。
  \item 命令参数中可能包含空格。
\end{itemize}

\subsection{逐字符扫描模板}

\begin{lstlisting}
string s;
cin >> s;

for (int i = 0; i < (int)s.size(); i++) {
    char c = s[i];
    if (isdigit((unsigned char)c)) {
        int j = i;
        while (j < (int)s.size() && isdigit((unsigned char)s[j])) j++;
        string number = s.substr(i, j - i);
        i = j - 1;
    } else if (isalpha((unsigned char)c)) {
        int j = i;
        while (j < (int)s.size() && isalpha((unsigned char)s[j])) j++;
        string word = s.substr(i, j - i);
        i = j - 1;
    } else {
        // symbol
    }
}
\end{lstlisting}

\subsection{解析 key=value}

\begin{lstlisting}
string s;
cin >> s;
int pos = (int)s.find('=');
string key = s.substr(0, pos);
string value = s.substr(pos + 1);
\end{lstlisting}

\subsection{路径解析}

\begin{lstlisting}
vector<string> parsePath(const string& path) {
    vector<string> parts;
    string cur;
    for (char c : path) {
        if (c == '/') {
            if (!cur.empty()) {
                parts.push_back(cur);
                cur.clear();
            }
        } else {
            cur.push_back(c);
        }
    }
    if (!cur.empty()) parts.push_back(cur);
    return parts;
}
\end{lstlisting}

\subsection{括号匹配}

\begin{lstlisting}
bool checkBrackets(const string& s) {
    stack<char> st;
    for (char c : s) {
        if (c == '(' || c == '[' || c == '{') {
            st.push(c);
        } else if (c == ')' || c == ']' || c == '}') {
            if (st.empty()) return false;
            char t = st.top();
            st.pop();
            if (c == ')' && t != '(') return false;
            if (c == ']' && t != '[') return false;
            if (c == '}' && t != '{') return false;
        }
    }
    return st.empty();
}
\end{lstlisting}

\subsection{字符串解析易错点}

\begin{itemize}
  \item \Code{cin >> s} 不能读取空格。
  \item \Code{getline} 前可能需要 \Code{cin.ignore()}。
  \item \Code{find} 找不到时返回 \Code{string::npos}。
  \item \Code{substr(pos, len)} 的第二个参数是长度，不是结束位置。
  \item 连续分隔符和开头/结尾分隔符要单独考虑。
  \item 字符串数字可能超过 \Code{long long}。
\end{itemize}

\section{表达式解析入口}\label{sec:07-22}

\subsection{什么题会用到}

题面出现：

\begin{itemize}
  \item 表达式、括号、运算符。
  \item 加减乘除优先级。
  \item 逻辑表达式。
  \item 嵌套结构。
\end{itemize}

\subsection{只含加减的表达式}

\begin{lstlisting}
long long evalPlusMinus(const string& s) {
    long long ans = 0;
    long long num = 0;
    int sign = 1;

    for (int i = 0; i <= (int)s.size(); i++) {
        if (i < (int)s.size() && isspace((unsigned char)s[i])) continue;
        if (i < (int)s.size() && isdigit((unsigned char)s[i])) {
            num = num * 10 + (s[i] - '0');
        } else {
            ans += sign * num;
            num = 0;
            if (i < (int)s.size()) {
                if (s[i] == '+') sign = 1;
                else if (s[i] == '-') sign = -1;
            }
        }
    }
    return ans;
}
\end{lstlisting}

完整整数四则运算见 \Tref{610}；下面仅展示读取一个数字的基础片段。

\subsection{递归解析思想}

复杂表达式通常需要递归下降。新手先掌握思想：

\begin{lstlisting}
string s;
int pos = 0;

long long parseNumber() {
    long long x = 0;
    while (pos < (int)s.size() && isdigit((unsigned char)s[pos])) {
        x = x * 10 + (s[pos] - '0');
        pos++;
    }
    return x;
}
\end{lstlisting}

\Warn{表达式题如果没练过，不要强行写复杂递归下降。}先拿没有括号、没有优先级的部分分。

\section{缓存和队列模拟}\label{sec:07-23}

\subsection{什么题会用到}

题面出现：

\begin{itemize}
  \item 缓存、页面置换、最近使用。
  \item 队列、排队、服务窗口。
  \item 先进先出、最近最少使用。
  \item 容量限制，满了要删除一个。
\end{itemize}

\subsection{FIFO 缓存}

\begin{lstlisting}
int cap;
cin >> cap;
queue<int> q;
set<int> inCache;

for (int x : requests) {
    if (cap <= 0) continue;
    if (inCache.count(x)) {
        continue;
    }
    if ((int)q.size() == cap) {
        int old = q.front();
        q.pop();
        inCache.erase(old);
    }
    q.push(x);
    inCache.insert(x);
}
\end{lstlisting}

\subsection{LRU 缓存入口}

LRU 可用 \Code{list + unordered\_map}：链表记访问顺序，映射保存每个对象的位置。

\begin{lstlisting}
int cap;
list<int> order; // front is most recent
unordered_map<int, list<int>::iterator> pos;

void access(int x) {
    if (cap <= 0) return;
    if (pos.count(x)) {
        order.erase(pos[x]);
    } else if ((int)order.size() == cap) {
        int old = order.back();
        order.pop_back();
        pos.erase(old);
    }
    order.push_front(x);
    pos[x] = order.begin();
}
\end{lstlisting}

\subsection{队列模拟易错点}

\begin{itemize}
  \item 队列为空时不能取 \Code{front}。
  \item 缓存容量可能为 0，要看题目是否允许。
  \item 命中缓存时是否需要更新时间。
  \item 删除元素时，辅助集合或 map 也要同步删除。
\end{itemize}

\section{懒删除优先队列}\label{sec:07-24}

\subsection{什么题会用到}

题面出现：

\begin{itemize}
  \item 需要动态删除任务。
  \item 还要每次取最大/最小。
  \item \Code{priority\_queue} 中间删除不方便。
\end{itemize}

\subsection{懒删除思想}

不直接从堆里删除旧元素，而是在取堆顶时检查它是否还有效。

\begin{lstlisting}
struct Task {
    int priority;
    int id;
};

struct Cmp {
    bool operator()(const Task& a, const Task& b) const {
        if (a.priority != b.priority) return a.priority < b.priority;
        return a.id > b.id;
    }
};

priority_queue<Task, vector<Task>, Cmp> pq;
map<int, int> currentPriority;

void addOrUpdate(int id, int priority) {
    currentPriority[id] = priority;
    pq.push({priority, id});
}

void removeTask(int id) {
    currentPriority.erase(id);
}

Task getTop() {
    while (!pq.empty()) {
        Task t = pq.top();
        if (currentPriority.count(t.id) &&
            currentPriority[t.id] == t.priority) {
            return t;
        }
        pq.pop();
    }
    return {-1, -1};
}
\end{lstlisting}

\Warn{懒删除会让堆里有废数据。}每次取 \Code{top} 前必须清理过期元素。

\section{时间和日程模拟}\label{sec:07-25}

\subsection{什么题会用到}

题面出现：

\begin{itemize}
  \item 时间段、会议、日程。
  \item 判断冲突。
  \item 开始时间、结束时间。
  \item 按时间排序处理。
\end{itemize}

\subsection{时间转分钟}

\begin{lstlisting}
int toMinute(const string& s) {
    int h = stoi(s.substr(0, 2));
    int m = stoi(s.substr(3, 2));
    return h * 60 + m;
}
\end{lstlisting}

\subsection{判断区间冲突}

两个半开区间 $[l1,r1)$ 和 $[l2,r2)$ 冲突：

\begin{lstlisting}
bool overlap(int l1, int r1, int l2, int r2) {
    return max(l1, l2) < min(r1, r2);
}
\end{lstlisting}

如果题目是闭区间，条件要改。

\subsection{扫描事件统计同时进行数量}

\begin{lstlisting}
vector<pair<int, int>> events;
for (auto [l, r] : intervals) {
    if (l >= r) continue; // empty half-open interval
    events.push_back({l, 1});
    events.push_back({r, -1});
}

sort(events.begin(), events.end(), [](auto a, auto b) {
    if (a.first != b.first) return a.first < b.first;
    return a.second < b.second; // end before start for [l,r)
});

int cur = 0, best = 0;
for (auto [time, delta] : events) {
    cur += delta;
    best = max(best, cur);
}
\end{lstlisting}

\subsection{时间题易错点}

\begin{itemize}
  \item 时间段是闭区间还是半开区间。
  \item 结束时间和开始时间相同是否冲突。
  \item 跨天如何处理。
  \item 同一时间开始和结束谁先处理。
\end{itemize}

\section{树形结构模拟}\label{sec:07-26}

\subsection{什么题会用到}

题面出现：

\begin{itemize}
  \item 层级、父子、目录、组织结构。
  \item 每个节点有父节点。
  \item 查询祖先、子树、路径。
\end{itemize}

\subsection{父子关系存储}

\begin{lstlisting}
int n;
cin >> n;
vector<vector<int>> children(n + 1);
vector<int> parent(n + 1, 0);

for (int i = 2; i <= n; i++) {
    int p;
    cin >> p;
    parent[i] = p;
    children[p].push_back(i);
}
\end{lstlisting}

\subsection{DFS 遍历树}

\begin{lstlisting}
vector<int> depth(n + 1, 0);

void dfsTree(int u) {
    for (int v : children[u]) {
        depth[v] = depth[u] + 1;
        dfsTree(v);
    }
}
\end{lstlisting}

\subsection{子树大小}

\begin{lstlisting}
vector<int> subtree(n + 1, 1);

void calcSubtree(int u) {
    subtree[u] = 1;
    for (int v : children[u]) {
        calcSubtree(v);
        subtree[u] += subtree[v];
    }
}
\end{lstlisting}

\subsection{树题易错点}

\begin{itemize}
  \item 根节点是谁。
  \item 输入给的是父节点，还是边。
  \item 节点编号从 0 还是 1 开始。
  \item 树很深时递归可能爆栈。
  \item 如果需要频繁祖先查询，简单爬父亲可能超时。
\end{itemize}

\section{权限和规则匹配}\label{sec:07-27}

\subsection{什么题会用到}

题面出现：

\begin{itemize}
  \item 用户、角色、权限。
  \item 规则匹配。
  \item 黑名单、白名单。
  \item 多条规则按顺序匹配，第一条生效。
  \item 默认允许或默认拒绝。
\end{itemize}

\subsection{规则结构体}

\begin{lstlisting}
struct Rule {
    string user;
    string action;
    string resource;
    bool allow;
};

vector<Rule> rules;
\end{lstlisting}

\subsection{规则匹配函数}

\begin{lstlisting}
bool match(const Rule& r, const string& user,
           const string& action, const string& resource) {
    if (r.user != "*" && r.user != user) return false;
    if (r.action != "*" && r.action != action) return false;
    if (r.resource != "*" && r.resource != resource) return false;
    return true;
}
\end{lstlisting}

\subsection{按顺序匹配}

\begin{lstlisting}
bool checkPermission(const string& user,
                     const string& action,
                     const string& resource) {
    for (const Rule& r : rules) {
        if (match(r, user, action, resource)) {
            return r.allow;
        }
    }
    return false; // default deny
}
\end{lstlisting}

\subsection{规则题易错点}

\begin{itemize}
  \item 多条规则是第一条生效，还是最后一条生效。
  \item 默认值是允许还是拒绝。
  \item 通配符如何匹配。
  \item 大小写是否敏感。
  \item 资源路径是否需要前缀匹配。
\end{itemize}

\section{网格和图混合题}\label{sec:07-28}

\subsection{什么题会用到}

题面出现：

\begin{itemize}
  \item 地图、障碍、传送门、钥匙、门。
  \item 不只是简单最短路，还有状态。
  \item 每个格子可能有属性。
  \item 移动过程中会改变状态。
\end{itemize}

\subsection{普通网格 BFS}

如果只有障碍和空地，用 B 题 BFS 模板即可。

\subsection{带状态 BFS 入口}

如果移动时还带着状态，例如是否拿到钥匙，可以把状态加入距离数组。

\begin{lstlisting}
struct Node {
    int x;
    int y;
    int state;
};

vector<vector<vector<int>>> dist(n,
    vector<vector<int>>(m, vector<int>(S, -1)));

queue<Node> q;
dist[sx][sy][0] = 0;
q.push({sx, sy, 0});

while (!q.empty()) {
    Node cur = q.front();
    q.pop();

    for (int d = 0; d < 4; d++) {
        int nx = cur.x + dx[d];
        int ny = cur.y + dy[d];
        int ns = cur.state;
        if (nx < 0 || nx >= n || ny < 0 || ny >= m) continue;
        if (g[nx][ny] == '#') continue;
        // Check door permission, then update ns from g[nx][ny].
        // Do not index the grid before the bounds check.
        if (dist[nx][ny][ns] != -1) continue;
        dist[nx][ny][ns] = dist[cur.x][cur.y][cur.state] + 1;
        q.push({nx, ny, ns});
    }
}
\end{lstlisting}

本段仅是状态扩展框架，必须按题意补全门禁检查、状态转移和起点状态。空间约为 $4nmS$ 字节，先算总量。

\Warn{状态数不能太大。}如果状态有 $2^k$，先看 $k$ 是否很小。$k \le 10$ 常见可做，$k=30$ 通常不行。

\subsection{图模拟策略}

如果题目是图但规则很多：

\begin{itemize}
  \item 先把图建对：有向/无向、权值、编号。
  \item 再处理规则：哪些边能走，什么时候能走。
  \item 如果边权全为 1，用 BFS。
  \item 如果边权为正，用 Dijkstra。
  \item 如果只是连通关系，用并查集或 DFS。
\end{itemize}

\section{大模拟易错点}\label{sec:07-29}

\begin{itemize}
  \item 同一时间多个事件的顺序。
  \item 相同优先级的并列规则。
  \item 删除元素后迭代器失效。
  \item 查询不存在对象。
  \item 重复添加同一对象。
  \item 空集合时取最大、最小。
  \item 最后一轮事件是否处理。
  \item 题目中的编号是 0 开始还是 1 开始。
  \item 输出顺序是否要求按编号、时间或字典序。
\end{itemize}

\section{C 题分阶段实现}\label{sec:07-30}

\subsection{不要一次写完}

每写完一种命令就运行一次，检查前后状态。可以按下面的顺序逐步完成：

\begin{enumerate}
  \item 只写读入，打印读到的数据检查。
  \item 写数据结构，不写复杂规则。
  \item 实现第一种命令或最基础操作。
  \item 跑样例，看是否能得到部分中间结果。
  \item 再补第二种命令。
  \item 最后处理特殊规则、并列规则、异常情况。
\end{enumerate}

\subsection{命令题实现顺序}

\begin{longtable}{p{3cm}p{6cm}p{5cm}}
\toprule
阶段 & 做什么 & 目标 \\
\midrule
\endfirsthead
\toprule
阶段 & 做什么 & 目标 \\
\midrule
\endhead
1 & 读入所有命令，只识别命令名 & 确认输入没错。 \\
2 & 实现新增/初始化类命令 & 建立基本状态。 \\
3 & 实现查询类命令 & 能输出最简单答案。 \\
4 & 实现修改/删除类命令 & 状态开始完整。 \\
5 & 实现特殊规则 & 冲满分。 \\
\bottomrule
\end{longtable}

\subsection{模拟题实现顺序}

\begin{enumerate}
  \item 先实现没有特殊情况的主流程。
  \item 再实现边界检查。
  \item 再实现冲突和优先级。
  \item 最后实现异常输入或题目特别说明的细节。
\end{enumerate}

\section{拆分函数时的检查}\label{sec:07-31}

代码较长时，把解析、查询和修改拆开，定位问题会方便一些。

\subsection{建议函数类型}

\begin{longtable}{p{4cm}p{5cm}p{5cm}}
\toprule
函数类型 & 示例 & 作用 \\
\midrule
\endfirsthead
\toprule
函数类型 & 示例 & 作用 \\
\midrule
\endhead
读入函数 & \Code{readInput()} & 避免读入和处理混在一起。 \\
解析函数 & \Code{parsePath(s)} & 字符串处理单独测试。 \\
查询函数 & \Code{query(id)} & 查询逻辑独立。 \\
修改函数 & \Code{update(id,x)} & 状态改变集中处理。 \\
检查函数 & \Code{canMove(x,y)} & 规则判断集中处理。 \\
输出函数 & \Code{printAnswer()} & 控制格式。 \\
\bottomrule
\end{longtable}

\subsection{函数设计原则}

\begin{itemize}
  \item 一个函数只做一类事情。
  \item 函数名直接反映题目操作。
  \item 一个函数尽量只处理一类操作；几种互不相关的规则可以拆开。
  \item 修改状态的函数要特别小心副作用。
  \item 查询函数尽量不要修改状态，除非题目要求。
\end{itemize}

\subsection{推荐代码顺序}

\begin{lstlisting}
// constants and types
struct ...

// global variables
int n, m;
vector<...> ...

// helper functions
bool inside(...);
vector<string> split(...);

// operation functions
void handleAdd(...);
void handleQuery(...);

// main solve
void solve() { ... }
\end{lstlisting}

\section{分层测试方法}\label{sec:07-32}

\subsection{第一层：解析测试}

只测试输入解析，不测算法。

\begin{itemize}
  \item 路径是否被正确分割。
  \item 命令参数是否读对。
  \item 数字字符串是否转换正确。
  \item 空格和换行是否处理正确。
\end{itemize}

\subsection{第二层：单操作测试}

每种操作单独测：

\begin{itemize}
  \item 只 add 一次。
  \item 只 query 一次。
  \item add 后 query。
  \item delete 后 query。
\end{itemize}

\subsection{第三层：组合测试}

测试多操作交替：

\begin{itemize}
  \item add、delete、add 同一个对象。
  \item 多个对象按优先级排序。
  \item 同一时间多个事件。
  \item 状态反复切换。
\end{itemize}

\subsection{第四层：边界测试}

测试异常和边界：

\begin{itemize}
  \item 空集合查询。
  \item 不存在对象。
  \item 重复对象。
  \item 最大数据附近。
  \item 最后一条命令。
\end{itemize}

\section{C 题常见失分原因}\label{sec:07-33}

\begin{longtable}{p{4cm}p{5cm}p{5cm}}
\toprule
失分原因 & 表现 & 预防 \\
\midrule
\endfirsthead
\toprule
失分原因 & 表现 & 预防 \\
\midrule
\endhead
读题漏规则 & 样例某一步对不上 & 三遍读题，列规则表。 \\
状态没初始化 & 第一次查询就错 & 结构体构造或统一 init。 \\
命令读错 & 后续全错 & 先调试输出 op 和参数。 \\
比较器错 & 排名/优先级错 & 手写 3 个对象测试 top/order。 \\
删除不同步 & 查到已删除对象 & 容器、set、map 同步维护。 \\
最后事件漏处理 & 前面都对，最后错 & 检查循环结束后的收尾。 \\
输出顺序错 & 数值对但顺序错 & 看题目要求按什么顺序。 \\
复杂度过高 & 小样例过，大数据超时 & 估算每条命令复杂度。 \\
\bottomrule
\end{longtable}

\section{C 题本地调试}\label{sec:07-34}

\subsection{调试输出建议}

调试时输出状态：

\begin{lstlisting}
void debugTask(const Task& t) {
    cerr << "id=" << t.id
         << " status=" << t.status
         << " priority=" << t.priority << '\n';
}
\end{lstlisting}

处理命令时：

\begin{lstlisting}
cerr << "op=" << op << '\n';
\end{lstlisting}

\Warn{提交前删除所有调试输出。}C 题代码长，尤其容易漏掉 \Code{cerr} 或临时 \Code{cout}。

\subsection{缩小样例}

如果官方样例很长，自己造小样例：

\begin{itemize}
  \item 1 个对象，1 条命令。
  \item 1 个对象，2 条命令。
  \item 2 个对象，测试排序或优先级。
  \item 不存在对象的查询。
  \item 重复添加和删除。
\end{itemize}

\subsection{定位 bug 的方法}

\begin{enumerate}
  \item 先确认读入正确。
  \item 再确认每条命令被识别正确。
  \item 再确认命令前状态正确。
  \item 再确认命令后状态正确。
  \item 最后确认输出格式。
\end{enumerate}

\section{C 题常见边界库}\label{sec:07-35}

\begin{longtable}{p{4cm}p{5cm}p{5cm}}
\toprule
题型 & 边界 & 检查目标 \\
\midrule
\endfirsthead
\toprule
题型 & 边界 & 检查目标 \\
\midrule
\endhead
命令模拟 & 只有 1 条命令、没有查询、连续查询 & 初始化和输出。 \\
新增删除 & 重复新增、删除不存在、删除后查询 & 对象生命周期。 \\
状态机 & 非法状态转换、终止状态继续操作 & 状态规则。 \\
事件模拟 & 同一时间多个事件、开始和结束同一时刻 & 事件排序。 \\
优先级 & 优先级相同、时间相同、编号相同附近 & 比较器。 \\
字符串解析 & 空字段、连续分隔符、路径以分隔符结尾 & split 和 substr。 \\
网格 & 起点等于终点、无路、只有一个格子 & BFS 初始化。 \\
带状态 BFS & 拿不到状态、重复状态、状态数最大 & dist 维度。 \\
图 & 自环、重边、不连通、有向/无向 & 建图。 \\
\bottomrule
\end{longtable}

\section{C 题部分分策略}\label{sec:07-36}
先看子任务限制。比如“没有删除命令”“没有嵌套结构”或“对象数量很少”，都可能允许更短的实现。选定一档后，把这一档会出现的规则写完整；只实现几条命令，却忽略它们会遇到的并列和边界规则，未必能得分。

\section{按子任务逐步完成}\label{sec:07-37}
可以按下面的顺序写，每完成一步就用小样例检查一次。

\subsection{读入与基本操作}
确认所有输入都能读完，每种命令的参数个数正确。实现初始状态和最简单的操作，检查一次新增或修改之后，查询能否得到预期结果。

\subsection{完成选定子任务}
补齐这个子任务会用到的所有命令、并列规则和异常处理。再试重复添加、删除不存在对象、空集合查询等边界。测试点有多次输出时，所有输出都要符合题意。

\subsection{扩大适用范围}
时间允许时，再加入嵌套、更复杂的规则或更快的数据结构。修改前保留已经检查过的版本，新版本通过对照后再替换。若题目采用捆绑测试，还要满足该子任务的全部要求。

\section{来不及写完时}\label{sec:07-38}

\subsection{写不完大模拟}

\begin{itemize}
  \item 保留已完成命令，不要为了补复杂规则改崩基础逻辑。
  \item 先完成子任务所需的状态修改，再检查相应查询输出。
  \item 特殊规则用清晰的 \Code{if} 补，不要重构大段代码。
\end{itemize}

\subsection{字符串解析写不出来}

\begin{itemize}
  \item 若有不含空格、不含嵌套的子任务，可以先处理这一档。
  \item 先用 \Code{split} 处理简单分隔符。
  \item 没有嵌套的子任务仍要正确处理其中所有 token 和输出。
\end{itemize}

\subsection{状态 BFS 写不出来}

\begin{itemize}
  \item 若子任务不含特殊状态，可以先写普通 BFS。
  \item 如果状态只有 2 种，尝试 \Code{dist[x][y][2]}。
  \item 如果有不含特殊状态的子任务，可以先完成这一档。
\end{itemize}

\subsection{复杂度可能超时}

\begin{itemize}
  \item 先估每条命令复杂度。
  \item 能用 map 查找就不要全表扫描。
  \item 能预处理就不要每次重复计算。
  \item 时间不够时，保留暴力拿小数据分。
\end{itemize}

\section{不同题型的实现顺序}\label{sec:07-39}

\subsection{如果是大模拟}

\begin{enumerate}
  \item 先读懂输入输出。
  \item 建结构体保存对象。
  \item 实现最简单命令。
  \item 实现查询输出。
  \item 再补删除、修改、异常情况。
  \item 最后补优先级、并列规则。
\end{enumerate}

\subsection{如果是字符串解析}

\begin{enumerate}
  \item 先确认用 \Code{cin} 还是 \Code{getline}。
  \item 先把字符串拆成 token。
  \item 能用 \Code{split} 就不要写复杂解析。
  \item 如果有括号，先写括号匹配。
  \item 如果有嵌套，先拿无嵌套部分分。
\end{enumerate}

\subsection{如果是图或网格}

\begin{enumerate}
  \item 先把图或网格读入正确。
  \item 判断是最短路、连通块还是模拟移动。
  \item 无权最短路用 BFS。
  \item 带状态时先估算状态数量。
  \item 只有子任务确实不含特殊状态时，才能用普通 BFS 简化。
\end{enumerate}

\subsection{如果是规则匹配}

\begin{enumerate}
  \item 把每条规则保存为结构体。
  \item 单独写 \Code{match} 函数。
  \item 确认第一条生效还是最后一条生效。
  \item 确认默认值。
  \item 用 2 到 3 条规则手测。
\end{enumerate}

\section{C 题模板索引}\label{sec:07-40}

\begin{longtable}{p{4cm}p{5cm}p{5cm}}
\toprule
题面关键词 & 使用模板 & 本章位置 \\
\midrule
\endfirsthead
\toprule
题面关键词 & 使用模板 & 本章位置 \\
\midrule
\endhead
多命令 & 命令分发 & 命令解析 \\
对象状态 & struct + 状态机 & 状态设计、状态机模拟 \\
按时间发生 & Event + sort & 事件模拟 \\
优先级 & priority\_queue & 优先级模拟 \\
排名 & 多关键字 sort & 复杂排序和排名 \\
表格、矩阵 & vector<vector<int>> & 表格和矩阵模拟 \\
路径、目录 & split + normalizePath & 目录树和路径模拟 \\
URL、参数 & parseUrl & URL 路由和参数解析 \\
配置、键值 & key=value & 配置和 JSON 类解析 \\
括号、表达式 & stack / 递归解析 & 字符串解析、表达式解析 \\
缓存 & queue / list + map & 缓存和队列模拟 \\
动态删除最大值 & 懒删除堆 & 懒删除优先队列 \\
时间段冲突 & overlap / events & 时间和日程模拟 \\
父子层级 & children + DFS & 树形结构模拟 \\
地图状态 & state BFS & 网格和图混合题 \\
\bottomrule
\end{longtable}

\section{C 题考场速查页}\label{sec:07-41}

\subsection{读题 5 分钟内要确定}

\begin{itemize}
  \item 是大模拟、字符串解析、图网格、数据结构，还是 DP。
  \item 有哪些对象。
  \item 每个对象有哪些状态。
  \item 有哪些命令或事件。
  \item 输出发生在每次查询还是最后。
  \item 数据范围是否允许全表扫描。
\end{itemize}

\subsection{写代码前先想清楚}

\begin{itemize}
  \item \Code{struct} 里有哪些字段。
  \item 用 \Code{vector}、\Code{map}、\Code{set} 还是堆。
  \item 每个命令对应哪个函数。
  \item 比较器的排序规则。
  \item 不存在、重复、越界怎么处理。
\end{itemize}

\subsection{最常用代码形状}

\begin{lstlisting}
struct Obj {
    int id;
    int status;
};

map<int, Obj> obj;

void handleAdd() {}
void handleDelete() {}
void handleQuery() {}

void solve() {
    int q;
    cin >> q;
    while (q--) {
        string op;
        cin >> op;
        if (op == "add") handleAdd();
        else if (op == "delete") handleDelete();
        else if (op == "query") handleQuery();
    }
}
\end{lstlisting}

\subsection{最后 10 分钟检查}

\begin{itemize}
  \item 删除调试输出。
  \item 检查所有命令是否都处理。
  \item 检查输出大小写、空格、换行。
  \item 检查比较器并列规则。
  \item 检查空集合、找不到对象。
  \item 检查最后一个事件。
  \item 提交一个当前最稳版本。
\end{itemize}

\section{C 题提交前检查}\label{sec:07-42}

\begin{itemize}
  \item 是否删除所有调试输出。
  \item 每种命令是否都处理了。
  \item 每种状态是否都有初始值。
  \item 不存在对象时是否按题意处理。
  \item 重复对象是否按题意处理。
  \item 排序和优先级比较器是否正确。
  \item 字符串读取是否会漏掉空格。
  \item 多组数据是否清空所有容器。
  \item 最后一个事件或最后一段数据是否处理。
  \item 输出顺序是否符合题目。
\end{itemize}

\section{C 题模板练习清单}\label{sec:07-43}

\begin{longtable}{p{4cm}p{4cm}p{6cm}}
\toprule
模板 & 熟练要求 & 用途 \\
\midrule
\endfirsthead
\toprule
模板 & 熟练要求 & 用途 \\
\midrule
\endhead
结构体 + map 保存对象 & 建议掌握 & 对象状态模拟。 \\
命令分发 & 建议掌握 & 多命令题。 \\
getline + stringstream & 建议掌握 & 整行命令。 \\
split & 照着能写 & 路径、配置解析。 \\
状态机 & 照着能写 & 多状态对象。 \\
事件排序 & 照着能写 & 时间顺序模拟。 \\
priority\_queue 自定义比较 & 练过再用 & 优先级模拟。 \\
括号匹配 & 照着能写 & 表达式/嵌套结构。 \\
带状态 BFS & 理解后再用 & 网格状态题。 \\
\bottomrule
\end{longtable}
